\scanpage{84}
NB. Accouplons $\mathbf Z^r$ et $\mathbf Z^r$ par
\[
((n_i),(n'_i))=\sum\delta_i n_i n'_i
\]
Alors $\lambda$ et $\mu$ sont formellement transposés par cet accouplement,
\deleted{\unsure{parfait}}, donc $\operatorname{Ker}\mu\simeq E^*(s)$ est accouplé à $\operatorname{Coker}\lambda\simeq E(s)$, par la formule
\[
(\alpha(x),y^*)=(x,\beta y^*)\qquad x\in\deleted{\unsure{E(s)|E_{(s)}}}\inserted{\mathbf Z^r},\ y^*\in E^*(s).
\]
\[
E(s)\times E^*(s)\longrightarrow\mathbf Z
\]
\deleted{Par suite,}\inserted{Je dis que} l'on a aussi
\[
(x,u(y))=(y,u(x))\qquad\text{pour }x,y\in E(s)
\]
En effet, on peut supposer que
\[
x=\alpha(\underbrace{(n_i)}_{\underline n}),\quad y=\alpha(\underbrace{(n'_j)}_{\underline n'}),
\]
et les deux mbs sont resp.
\[
\begin{aligned}
(\alpha(\underline n),u(y))&=(\underline n,\beta u\alpha(\underline n'))\\
(\alpha(\underline n'),u(x))&=(\underline n',\beta u\alpha(\underline n))
\end{aligned}
\]
Or on a
\[
(\underline n,\beta u\alpha(\underline n'))=\sum n_i n'_j(C_i C_j)
\]
qui est bien symétrique en \deleted{\unsure{$\in$|et}} $\underline n,\underline n'$.

Si les $\delta_i$ sont tous $1$, \inserted{(p.ex $k(s)$ p-fait)} alors on trouve \unsure{pr|pour|par} l'accouplement précédent
\[
E^*(s)\simeq\operatorname{Hom}(E(s),\mathbf Z)
\]
et la relation symétrique si on a de plus $d=1$ i.e. $d^t=1$.
Dans ce cas cette dualité, et le fait que $u:E(s)\unsure{\relbar\joinrel\relbar|\longrightarrow}E^*(s)$ \unsure{:} est autoadjoint, donne pour le $\operatorname{Coker}u=\Psi(s)$ une \uline{autodualité} canonique (à valeurs dans $\mathbf Q/\mathbf Z$).
Cette autodualité \deleted{\unsure{est liée à}}\inserted{exprime} l'obstruction à étendre la biextension qu'on a de $(Q'_\eta\times Q'_\eta)=(P^0_\eta\times P^0_\eta)$ par $G_{m\eta}$, par une biextension de $Q',Q'$\margin{modèles de Néron} par $G_m$
\margin{83}
